%% Supplementary material for "EduLogGen: A Python package for generating and
%% validating synthetic educational interaction logs".
\documentclass[11pt]{article}
\usepackage[margin=2.2cm]{geometry}
\usepackage{amssymb}
\usepackage{booktabs}
\usepackage{graphicx}
\usepackage{hyperref}

\title{Supplementary material\\[0.3em]\large EduLogGen: A Python package for generating and validating synthetic educational interaction logs}
\author{Anis Bey}
\date{}

\begin{document}
\maketitle

\section*{S1. Protocols}

All analyses use seed 0 and EduLogGen v1.5.0. Scripts:
\texttt{studies/oulad/} and \texttt{studies/ednet/} (data preparation and
runners), \texttt{studies/study.py} (analyses), \texttt{studies/report.py}
(tables and figures).

\paragraph{Split} Learners are split 70/30 into training and holdout sets
(\texttt{session\_fidelity\_v1}). Every analysis fits on training learners
only.

\paragraph{Fidelity} Each generator is fitted on the training learners and
generates three samples (seeds derived from seed 0) with as many sessions as
the holdout; each sample is compared with the holdout. The real-versus-real
reference compares the training learners, treated as a sample, with the
holdout once per dataset; it is computed on the same holdout as the
generator samples but with the larger training set as the second sample.
Generators: \texttt{independent}; \texttt{markov} with order 1 and 2;
\texttt{semi\_markov} (order 1, empirical gap distributions, per-token
distribution if at least 5 gaps were observed, otherwise the pooled one).

\paragraph{Profiles} $k=4$ automatic profiles (features: number of sessions,
mean session length, log mean session duration, success rate, and the share
of each event type with at least 1\% of events, rarer types pooled;
standardised with training statistics, $z$-scores capped at $\pm3$;
k-means++ with 10 restarts). Stability: ARI between the training learners'
assignments under seed 0 and under seeds 1--9; ARI between the full-data
assignments and those of models fitted on ten random 80\% subsamples of the
training learners (all training learners assigned); holdout replication: ARI
between holdout learners assigned to the training profiles and to profiles
fitted on the holdout alone. Recovery: a synthetic population with as many
sessions as the training data, mixing the profiles in their observed shares,
is clustered again with the same $k$.

\paragraph{Memorisation} Synthetic data (Markov-2 and semi-Markov, fitted on
training learners, as many sessions as the holdout) and holdout learners are
each compared with the training learners. Session level: sessions with at
least three events identical to any training session, and to a training
session that occurs exactly once. Learner level: learners with at least two
sessions whose ordered sequence of sessions equals a training learner's; and
the Euclidean distance from each learner to the closest training learner in
standardised behavioural features (no cap), summarised by its median and the
share of learners at distance 0.

\paragraph{Detection} Unexpected transition: two events $(a, b)$ inserted at
a random position, where the pair $(a, b)$ never occurs among training
learners. Repetition: four events alternating between two event types of the
session (gap 1 day on OULAD, 5 s on EdNet). Each type is injected into a
disjoint random 5\% of sessions; a session carries at most one anomaly and
is labelled anomalous at session level. Insertions lengthen the session by
two or four events; the detector uses transitions only. Detector: mean
negative log-likelihood of a session's transitions (including the start) under
an add-one-smoothed bigram model fitted on training learners. Synthetic
background: the same injection into a semi-Markov sample of the holdout's
size. Artefact check: ROC-AUC of the detector score for separating normal
synthetic sessions (positive) from normal holdout sessions.

\section*{S2. Results per module and dataset}

Fidelity: mean $\pm$ SD over the three generation seeds.

\subsection*{OULAD (7 modules): AAA-2014J}
357 learners, 30{,}703 events, 9{,}914 sessions; wall time 46~s.

\begin{tabular}{l r r r r r r }
\toprule
Metric & Real vs real & Independent & Markov-1 & Markov-2 & Semi-Markov & GRU \\
\midrule
Event-type TVD $\downarrow$ & 0.025 & 0.022 $\pm$ 0.004 & 0.017 $\pm$ 0.004 & 0.010 $\pm$ 0.005 & 0.010 $\pm$ 0.004 & 0.040 $\pm$ 0.009 \\
Bigram TVD $\downarrow$ & 0.031 & 0.235 $\pm$ 0.008 & 0.050 $\pm$ 0.006 & 0.055 $\pm$ 0.007 & 0.057 $\pm$ 0.005 & 0.093 $\pm$ 0.011 \\
Transition JSD $\downarrow$ & 0.003 & 0.046 $\pm$ 0.003 & 0.004 $\pm$ 0.000 & 0.004 $\pm$ 0.000 & 0.004 $\pm$ 0.000 & 0.005 $\pm$ 0.001 \\
Session-length KS $\downarrow$ & 0.054 & 0.058 $\pm$ 0.004 & 0.051 $\pm$ 0.006 & 0.051 $\pm$ 0.006 & 0.058 $\pm$ 0.007 & 0.053 $\pm$ 0.011 \\
Inter-event time KS $\downarrow$ & 0.015 & 0.245 $\pm$ 0.000 & 0.245 $\pm$ 0.000 & 0.245 $\pm$ 0.000 & 0.007 $\pm$ 0.004 & 0.010 $\pm$ 0.002 \\
Top-10 path overlap $\uparrow$ & 0.900 & 0.700 $\pm$ 0.100 & 0.900 $\pm$ 0.000 & 0.833 $\pm$ 0.058 & 0.867 $\pm$ 0.058 & 0.833 $\pm$ 0.058 \\
\bottomrule
\end{tabular}

Profiles ($k=4$): seed ARI 0.91, 0.85, 0.94, 0.91, 0.95, 0.94, 0.95, 0.95, 0.92; subsample ARI 0.78, 0.82, 0.77, 0.87, 0.85, 0.90, 0.72, 0.89, 0.82, 0.71; holdout replication ARI 0.71; recovery ARI 0.91; NMI with final result (holdout) 0.116.

Memorisation (sessions with $\geq$3 events; matches against training learners): holdout: any 86.0\%, once-seen 7.1\%, trajectories 0.0\%, feature-identical 0.0\%, median distance 0.64; Markov-2: any 80.9\%, once-seen 8.0\%, trajectories 0.0\%, feature-identical 0.0\%, median distance 0.63; semi-Markov: any 77.5\%, once-seen 8.1\%, trajectories 0.0\%, feature-identical 0.0\%, median distance 0.64; GRU: any 81.3\%, once-seen 8.9\%, trajectories 0.0\%, feature-identical 0.0\%, median distance 0.70.

Detection (real background): repetition: ROC-AUC 0.83, AP 0.26; unexpected transition: ROC-AUC 0.98, AP 0.55; synthetic background ROC-AUC 0.92; synthetic-vs-real AUC 0.48.

\subsection*{OULAD (7 modules): BBB-2014J}
1{,}921 learners, 83{,}111 events, 36{,}982 sessions; wall time 188~s.

\begin{tabular}{l r r r r r r }
\toprule
Metric & Real vs real & Independent & Markov-1 & Markov-2 & Semi-Markov & GRU \\
\midrule
Event-type TVD $\downarrow$ & 0.030 & 0.031 $\pm$ 0.003 & 0.033 $\pm$ 0.002 & 0.037 $\pm$ 0.002 & 0.036 $\pm$ 0.001 & 0.055 $\pm$ 0.004 \\
Bigram TVD $\downarrow$ & 0.049 & 0.183 $\pm$ 0.003 & 0.078 $\pm$ 0.004 & 0.083 $\pm$ 0.003 & 0.078 $\pm$ 0.002 & 0.115 $\pm$ 0.004 \\
Transition JSD $\downarrow$ & 0.002 & 0.039 $\pm$ 0.003 & 0.003 $\pm$ 0.000 & 0.003 $\pm$ 0.000 & 0.003 $\pm$ 0.000 & 0.005 $\pm$ 0.000 \\
Session-length KS $\downarrow$ & 0.007 & 0.007 $\pm$ 0.003 & 0.008 $\pm$ 0.003 & 0.008 $\pm$ 0.003 & 0.009 $\pm$ 0.003 & 0.012 $\pm$ 0.000 \\
Inter-event time KS $\downarrow$ & 0.004 & 0.348 $\pm$ 0.000 & 0.348 $\pm$ 0.000 & 0.348 $\pm$ 0.000 & 0.004 $\pm$ 0.001 & 0.015 $\pm$ 0.004 \\
Top-10 path overlap $\uparrow$ & 1.000 & 0.900 $\pm$ 0.000 & 0.900 $\pm$ 0.000 & 0.900 $\pm$ 0.000 & 0.900 $\pm$ 0.000 & 0.900 $\pm$ 0.000 \\
\bottomrule
\end{tabular}

Profiles ($k=4$): seed ARI 1.00, 1.00, 1.00, 1.00, 1.00, 1.00, 1.00, 0.99, 1.00; subsample ARI 0.97, 0.45, 0.98, 0.63, 0.89, 0.96, 0.98, 0.93, 0.69, 0.98; holdout replication ARI 0.84; recovery ARI 0.86; NMI with final result (holdout) 0.141.

Memorisation (sessions with $\geq$3 events; matches against training learners): holdout: any 87.8\%, once-seen 5.3\%, trajectories 3.1\%, feature-identical 8.3\%, median distance 0.70; Markov-2: any 85.5\%, once-seen 5.1\%, trajectories 0.4\%, feature-identical 3.2\%, median distance 0.81; semi-Markov: any 85.9\%, once-seen 6.2\%, trajectories 0.9\%, feature-identical 2.4\%, median distance 0.78; GRU: any 86.3\%, once-seen 5.5\%, trajectories 1.1\%, feature-identical 3.0\%, median distance 0.81.

Detection (real background): repetition: ROC-AUC 0.87, AP 0.28; unexpected transition: ROC-AUC 0.97, AP 0.49; synthetic background ROC-AUC 0.91; synthetic-vs-real AUC 0.49.

\subsection*{OULAD (7 modules): CCC-2014J}
2{,}302 learners, 131{,}279 events, 47{,}579 sessions; wall time 264~s.

\begin{tabular}{l r r r r r r }
\toprule
Metric & Real vs real & Independent & Markov-1 & Markov-2 & Semi-Markov & GRU \\
\midrule
Event-type TVD $\downarrow$ & 0.023 & 0.022 $\pm$ 0.002 & 0.027 $\pm$ 0.003 & 0.033 $\pm$ 0.004 & 0.027 $\pm$ 0.000 & 0.057 $\pm$ 0.002 \\
Bigram TVD $\downarrow$ & 0.036 & 0.242 $\pm$ 0.002 & 0.076 $\pm$ 0.002 & 0.083 $\pm$ 0.002 & 0.078 $\pm$ 0.003 & 0.104 $\pm$ 0.008 \\
Transition JSD $\downarrow$ & 0.001 & 0.063 $\pm$ 0.001 & 0.002 $\pm$ 0.000 & 0.002 $\pm$ 0.000 & 0.001 $\pm$ 0.000 & 0.004 $\pm$ 0.000 \\
Session-length KS $\downarrow$ & 0.011 & 0.015 $\pm$ 0.000 & 0.013 $\pm$ 0.000 & 0.013 $\pm$ 0.000 & 0.013 $\pm$ 0.001 & 0.015 $\pm$ 0.003 \\
Inter-event time KS $\downarrow$ & 0.009 & 0.266 $\pm$ 0.000 & 0.266 $\pm$ 0.000 & 0.266 $\pm$ 0.000 & 0.020 $\pm$ 0.002 & 0.032 $\pm$ 0.001 \\
Top-10 path overlap $\uparrow$ & 0.800 & 0.600 $\pm$ 0.000 & 0.900 $\pm$ 0.000 & 0.733 $\pm$ 0.058 & 0.900 $\pm$ 0.000 & 0.700 $\pm$ 0.000 \\
\bottomrule
\end{tabular}

Profiles ($k=4$): seed ARI 1.00, 1.00, 1.00, 1.00, 1.00, 1.00, 0.99, 1.00, 1.00; subsample ARI 0.93, 0.70, 0.70, 0.68, 0.94, 0.89, 0.73, 0.87, 0.99, 0.75; holdout replication ARI 0.70; recovery ARI 0.78; NMI with final result (holdout) 0.103.

Memorisation (sessions with $\geq$3 events; matches against training learners): holdout: any 83.9\%, once-seen 6.2\%, trajectories 2.3\%, feature-identical 5.8\%, median distance 0.74; Markov-2: any 82.7\%, once-seen 5.8\%, trajectories 0.7\%, feature-identical 3.7\%, median distance 0.81; semi-Markov: any 80.3\%, once-seen 6.2\%, trajectories 0.6\%, feature-identical 1.9\%, median distance 0.79; GRU: any 84.6\%, once-seen 5.3\%, trajectories 0.4\%, feature-identical 1.0\%, median distance 0.75.

Detection (real background): repetition: ROC-AUC 0.82, AP 0.25; unexpected transition: ROC-AUC 0.98, AP 0.63; synthetic background ROC-AUC 0.89; synthetic-vs-real AUC 0.52.

\subsection*{OULAD (7 modules): DDD-2014J}
1{,}647 learners, 109{,}161 events, 38{,}264 sessions; wall time 259~s.

\begin{tabular}{l r r r r r r }
\toprule
Metric & Real vs real & Independent & Markov-1 & Markov-2 & Semi-Markov & GRU \\
\midrule
Event-type TVD $\downarrow$ & 0.017 & 0.017 $\pm$ 0.006 & 0.032 $\pm$ 0.001 & 0.043 $\pm$ 0.002 & 0.033 $\pm$ 0.002 & 0.069 $\pm$ 0.005 \\
Bigram TVD $\downarrow$ & 0.030 & 0.217 $\pm$ 0.004 & 0.085 $\pm$ 0.002 & 0.095 $\pm$ 0.000 & 0.086 $\pm$ 0.005 & 0.122 $\pm$ 0.007 \\
Transition JSD $\downarrow$ & 0.001 & 0.047 $\pm$ 0.002 & 0.002 $\pm$ 0.000 & 0.002 $\pm$ 0.000 & 0.002 $\pm$ 0.000 & 0.005 $\pm$ 0.000 \\
Session-length KS $\downarrow$ & 0.013 & 0.018 $\pm$ 0.000 & 0.015 $\pm$ 0.002 & 0.015 $\pm$ 0.002 & 0.012 $\pm$ 0.001 & 0.014 $\pm$ 0.002 \\
Inter-event time KS $\downarrow$ & 0.011 & 0.278 $\pm$ 0.000 & 0.278 $\pm$ 0.000 & 0.278 $\pm$ 0.000 & 0.022 $\pm$ 0.001 & 0.007 $\pm$ 0.004 \\
Top-10 path overlap $\uparrow$ & 1.000 & 0.800 $\pm$ 0.000 & 0.833 $\pm$ 0.058 & 1.000 $\pm$ 0.000 & 0.900 $\pm$ 0.000 & 1.000 $\pm$ 0.000 \\
\bottomrule
\end{tabular}

Profiles ($k=4$): seed ARI 1.00, 0.95, 1.00, 0.98, 0.96, 1.00, 1.00, 1.00, 1.00; subsample ARI 0.96, 0.97, 0.92, 0.74, 0.89, 0.94, 0.97, 0.93, 0.78, 0.98; holdout replication ARI 0.54; recovery ARI 0.84; NMI with final result (holdout) 0.145.

Memorisation (sessions with $\geq$3 events; matches against training learners): holdout: any 81.8\%, once-seen 7.2\%, trajectories 1.3\%, feature-identical 3.6\%, median distance 0.90; Markov-2: any 80.3\%, once-seen 6.4\%, trajectories 0.0\%, feature-identical 1.4\%, median distance 0.94; semi-Markov: any 76.8\%, once-seen 7.0\%, trajectories 0.0\%, feature-identical 2.4\%, median distance 0.90; GRU: any 79.4\%, once-seen 6.6\%, trajectories 0.0\%, feature-identical 1.5\%, median distance 0.92.

Detection (real background): repetition: ROC-AUC 0.79, AP 0.21; unexpected transition: ROC-AUC 0.97, AP 0.53; synthetic background ROC-AUC 0.88; synthetic-vs-real AUC 0.50.

\subsection*{OULAD (7 modules): EEE-2014J}
1{,}097 learners, 76{,}751 events, 26{,}198 sessions; wall time 137~s.

\begin{tabular}{l r r r r r r }
\toprule
Metric & Real vs real & Independent & Markov-1 & Markov-2 & Semi-Markov & GRU \\
\midrule
Event-type TVD $\downarrow$ & 0.019 & 0.019 $\pm$ 0.003 & 0.031 $\pm$ 0.003 & 0.039 $\pm$ 0.004 & 0.033 $\pm$ 0.009 & 0.036 $\pm$ 0.001 \\
Bigram TVD $\downarrow$ & 0.033 & 0.243 $\pm$ 0.006 & 0.069 $\pm$ 0.005 & 0.075 $\pm$ 0.007 & 0.073 $\pm$ 0.011 & 0.080 $\pm$ 0.001 \\
Transition JSD $\downarrow$ & 0.002 & 0.062 $\pm$ 0.002 & 0.002 $\pm$ 0.000 & 0.002 $\pm$ 0.001 & 0.002 $\pm$ 0.000 & 0.003 $\pm$ 0.000 \\
Session-length KS $\downarrow$ & 0.021 & 0.023 $\pm$ 0.001 & 0.023 $\pm$ 0.009 & 0.023 $\pm$ 0.009 & 0.021 $\pm$ 0.003 & 0.021 $\pm$ 0.006 \\
Inter-event time KS $\downarrow$ & 0.008 & 0.258 $\pm$ 0.000 & 0.258 $\pm$ 0.000 & 0.258 $\pm$ 0.000 & 0.003 $\pm$ 0.001 & 0.027 $\pm$ 0.005 \\
Top-10 path overlap $\uparrow$ & 0.900 & 0.600 $\pm$ 0.000 & 0.767 $\pm$ 0.058 & 0.700 $\pm$ 0.000 & 0.767 $\pm$ 0.058 & 0.700 $\pm$ 0.000 \\
\bottomrule
\end{tabular}

Profiles ($k=4$): seed ARI 0.97, 0.95, 0.94, 0.97, 0.97, 0.94, 0.89, 0.97, 0.99; subsample ARI 0.93, 0.94, 0.95, 0.92, 0.95, 0.97, 0.96, 0.95, 0.89, 0.94; holdout replication ARI 0.81; recovery ARI 0.88; NMI with final result (holdout) 0.302.

Memorisation (sessions with $\geq$3 events; matches against training learners): holdout: any 84.1\%, once-seen 6.8\%, trajectories 1.9\%, feature-identical 4.8\%, median distance 0.69; Markov-2: any 80.8\%, once-seen 6.7\%, trajectories 0.3\%, feature-identical 2.0\%, median distance 0.75; semi-Markov: any 78.3\%, once-seen 7.3\%, trajectories 0.0\%, feature-identical 1.4\%, median distance 0.81; GRU: any 82.0\%, once-seen 7.1\%, trajectories 0.3\%, feature-identical 1.4\%, median distance 0.77.

Detection (real background): repetition: ROC-AUC 0.83, AP 0.24; unexpected transition: ROC-AUC 0.97, AP 0.50; synthetic background ROC-AUC 0.90; synthetic-vs-real AUC 0.50.

\subsection*{OULAD (7 modules): FFF-2014J}
2{,}121 learners, 167{,}164 events, 50{,}431 sessions; wall time 316~s.

\begin{tabular}{l r r r r r r }
\toprule
Metric & Real vs real & Independent & Markov-1 & Markov-2 & Semi-Markov & GRU \\
\midrule
Event-type TVD $\downarrow$ & 0.022 & 0.022 $\pm$ 0.003 & 0.028 $\pm$ 0.002 & 0.030 $\pm$ 0.002 & 0.026 $\pm$ 0.001 & 0.027 $\pm$ 0.005 \\
Bigram TVD $\downarrow$ & 0.034 & 0.219 $\pm$ 0.004 & 0.061 $\pm$ 0.004 & 0.062 $\pm$ 0.003 & 0.060 $\pm$ 0.003 & 0.060 $\pm$ 0.005 \\
Transition JSD $\downarrow$ & 0.001 & 0.040 $\pm$ 0.002 & 0.002 $\pm$ 0.000 & 0.002 $\pm$ 0.000 & 0.002 $\pm$ 0.000 & 0.002 $\pm$ 0.000 \\
Session-length KS $\downarrow$ & 0.015 & 0.014 $\pm$ 0.004 & 0.018 $\pm$ 0.003 & 0.018 $\pm$ 0.003 & 0.013 $\pm$ 0.004 & 0.015 $\pm$ 0.003 \\
Inter-event time KS $\downarrow$ & 0.002 & 0.222 $\pm$ 0.000 & 0.222 $\pm$ 0.000 & 0.222 $\pm$ 0.000 & 0.003 $\pm$ 0.001 & 0.010 $\pm$ 0.003 \\
Top-10 path overlap $\uparrow$ & 0.900 & 0.700 $\pm$ 0.000 & 0.700 $\pm$ 0.000 & 0.833 $\pm$ 0.058 & 0.700 $\pm$ 0.000 & 0.900 $\pm$ 0.000 \\
\bottomrule
\end{tabular}

Profiles ($k=4$): seed ARI 0.93, 0.93, 0.92, 0.93, 0.94, 0.93, 0.93, 0.93, 0.98; subsample ARI 0.89, 0.90, 0.78, 0.64, 0.88, 0.96, 0.71, 0.91, 0.87, 0.95; holdout replication ARI 0.36; recovery ARI 0.92; NMI with final result (holdout) 0.199.

Memorisation (sessions with $\geq$3 events; matches against training learners): holdout: any 83.2\%, once-seen 5.8\%, trajectories 0.8\%, feature-identical 3.1\%, median distance 0.66; Markov-2: any 81.2\%, once-seen 6.7\%, trajectories 0.3\%, feature-identical 1.4\%, median distance 0.63; semi-Markov: any 79.3\%, once-seen 7.3\%, trajectories 0.0\%, feature-identical 1.8\%, median distance 0.64; GRU: any 81.9\%, once-seen 6.6\%, trajectories 0.5\%, feature-identical 1.4\%, median distance 0.62.

Detection (real background): repetition: ROC-AUC 0.80, AP 0.24; unexpected transition: ROC-AUC 0.97, AP 0.55; synthetic background ROC-AUC 0.89; synthetic-vs-real AUC 0.50.

\subsection*{OULAD (7 modules): GGG-2014J}
698 learners, 25{,}535 events, 12{,}467 sessions; wall time 54~s.

\begin{tabular}{l r r r r r r }
\toprule
Metric & Real vs real & Independent & Markov-1 & Markov-2 & Semi-Markov & GRU \\
\midrule
Event-type TVD $\downarrow$ & 0.017 & 0.019 $\pm$ 0.008 & 0.025 $\pm$ 0.004 & 0.028 $\pm$ 0.004 & 0.021 $\pm$ 0.006 & 0.042 $\pm$ 0.010 \\
Bigram TVD $\downarrow$ & 0.052 & 0.206 $\pm$ 0.011 & 0.096 $\pm$ 0.018 & 0.097 $\pm$ 0.017 & 0.082 $\pm$ 0.003 & 0.116 $\pm$ 0.011 \\
Transition JSD $\downarrow$ & 0.003 & 0.046 $\pm$ 0.005 & 0.006 $\pm$ 0.001 & 0.006 $\pm$ 0.002 & 0.006 $\pm$ 0.001 & 0.009 $\pm$ 0.002 \\
Session-length KS $\downarrow$ & 0.019 & 0.025 $\pm$ 0.005 & 0.014 $\pm$ 0.005 & 0.014 $\pm$ 0.005 & 0.018 $\pm$ 0.006 & 0.027 $\pm$ 0.003 \\
Inter-event time KS $\downarrow$ & 0.025 & 0.349 $\pm$ 0.000 & 0.349 $\pm$ 0.000 & 0.349 $\pm$ 0.000 & 0.030 $\pm$ 0.010 & 0.020 $\pm$ 0.008 \\
Top-10 path overlap $\uparrow$ & 0.800 & 0.600 $\pm$ 0.000 & 0.833 $\pm$ 0.058 & 0.700 $\pm$ 0.000 & 0.833 $\pm$ 0.058 & 0.700 $\pm$ 0.000 \\
\bottomrule
\end{tabular}

Profiles ($k=4$): seed ARI 0.97, 0.95, 0.95, 1.00, 0.96, 0.62, 0.97, 0.92, 0.97; subsample ARI 0.61, 0.90, 0.57, 0.90, 0.77, 0.62, 0.88, 0.78, 0.63, 0.95; holdout replication ARI 0.60; recovery ARI 0.85; NMI with final result (holdout) 0.198.

Memorisation (sessions with $\geq$3 events; matches against training learners): holdout: any 78.7\%, once-seen 8.5\%, trajectories 3.0\%, feature-identical 5.2\%, median distance 0.78; Markov-2: any 81.0\%, once-seen 7.2\%, trajectories 1.4\%, feature-identical 0.5\%, median distance 0.83; semi-Markov: any 76.9\%, once-seen 6.1\%, trajectories 0.5\%, feature-identical 0.5\%, median distance 0.80; GRU: any 78.0\%, once-seen 8.2\%, trajectories 0.5\%, feature-identical 2.8\%, median distance 0.80.

Detection (real background): repetition: ROC-AUC 0.83, AP 0.18; unexpected transition: ROC-AUC 0.97, AP 0.49; synthetic background ROC-AUC 0.91; synthetic-vs-real AUC 0.49.

\subsection*{EdNet-KT4: ednet-kt4-2000}
2{,}000 learners, 967{,}968 events, 10{,}580 sessions; wall time 4066~s.

\begin{tabular}{l r r r r r r }
\toprule
Metric & Real vs real & Independent & Markov-1 & Markov-2 & Semi-Markov & GRU \\
\midrule
Event-type TVD $\downarrow$ & 0.039 & 0.039 $\pm$ 0.000 & 0.047 $\pm$ 0.002 & 0.045 $\pm$ 0.002 & 0.048 $\pm$ 0.002 & 0.054 $\pm$ 0.005 \\
Bigram TVD $\downarrow$ & 0.052 & 0.655 $\pm$ 0.001 & 0.061 $\pm$ 0.002 & 0.059 $\pm$ 0.002 & 0.063 $\pm$ 0.002 & 0.094 $\pm$ 0.006 \\
Transition JSD $\downarrow$ & 0.004 & 0.488 $\pm$ 0.001 & 0.004 $\pm$ 0.000 & 0.004 $\pm$ 0.000 & 0.004 $\pm$ 0.000 & 0.009 $\pm$ 0.001 \\
Session-length KS $\downarrow$ & 0.084 & 0.079 $\pm$ 0.004 & 0.085 $\pm$ 0.003 & 0.085 $\pm$ 0.003 & 0.086 $\pm$ 0.003 & 0.089 $\pm$ 0.014 \\
Inter-event time KS $\downarrow$ & 0.019 & 0.507 $\pm$ 0.000 & 0.507 $\pm$ 0.000 & 0.507 $\pm$ 0.000 & 0.028 $\pm$ 0.000 & 0.032 $\pm$ 0.000 \\
Top-10 path overlap $\uparrow$ & 1.000 & 0.100 $\pm$ 0.000 & 0.800 $\pm$ 0.000 & 1.000 $\pm$ 0.000 & 0.800 $\pm$ 0.000 & 0.800 $\pm$ 0.000 \\
\bottomrule
\end{tabular}

Profiles ($k=4$): seed ARI 0.99, 0.99, 0.99, 0.99, 1.00, 0.99, 0.99, 1.00, 1.00; subsample ARI 0.99, 0.97, 0.99, 0.99, 0.99, 0.95, 0.97, 0.96, 0.96, 0.98; holdout replication ARI 0.93; recovery ARI 0.75.

Memorisation (sessions with $\geq$3 events; matches against training learners): holdout: any 20.2\%, once-seen 2.8\%, trajectories 0.0\%, feature-identical 0.0\%, median distance 0.42; Markov-2: any 2.9\%, once-seen 0.3\%, trajectories 0.0\%, feature-identical 0.0\%, median distance 1.63; semi-Markov: any 1.8\%, once-seen 0.2\%, trajectories 0.0\%, feature-identical 0.0\%, median distance 1.26; GRU: any 7.1\%, once-seen 1.4\%, trajectories 0.0\%, feature-identical 0.0\%, median distance 1.07.

Detection (real background): repetition: ROC-AUC 0.81, AP 0.21; unexpected transition: ROC-AUC 0.74, AP 0.17; synthetic background ROC-AUC 0.90; synthetic-vs-real AUC 0.49.



\end{document}
